\subsection{6. Peran Mahkamah Konstitusi}
Sejarah berdirinya Mahkamah Konstitusi (MK) diawali dengan diadopsinya ide Mahkamah Konstitusi (\textit{Constitutional Court}) dalam amendemen konstitusi yang dilakukan oleh Majelis Permusyawaratan Rakyat (MPR) pada tahun 2001 sebagaimana dirumuskan dalam ketentuan Pasal 24 ayat (2), Pasal 24C, dan Pasal 7B UUD NRI Tahun 1945. Kemudian DPR dan pemerintah menetapkan Undang-Undang Nomor 24 Tahun 2003 tentang Mahkamah Konstitusi pada tanggal 13 Agustus 2003.

Ada empat kewenangan dan satu kewajiban Mahkamah Konstitusi yang telah ditentukan dalam UUD NRI Tahun 1945 perubahan ketiga Pasal 24C ayat (1), yaitu sebagai berikut:
\begin{enumerate}
    \item Menguji (\textit{judicial review}) undang-undang terhadap UUD.
    \item Memutus sengketa kewenangan lembaga negara yang kewenangannya diberikan oleh UUD.
    \item Memutuskan pembubaran partai politik.
    \item Memutus perselisihan tentang hasil pemilihan umum.
    \item Selain itu, Mahkamah Konstitusi juga memiliki kewajiban memberikan putusan atas pendapat DPR mengenai dugaan pelanggaran hukum oleh Presiden dan/atau Wakil Presiden menurut UUD.
\end{enumerate}

Pengadilan yang dilakukan oleh Mahkamah Konstitusi merupakan pengadilan tingkat pertama dan terakhir yang putusannya bersifat final. Artinya, tidak ada upaya hukum lain atas putusan Mahkamah Konstitusi. Mahkamah Konstitusi memiliki peran sebagai berikut:
\begin{enumerate}
    \item Pengawal konstitusi (\textit{the guardian of constitution}).
    \item Penafsir akhir konstitusi (\textit{the final interpreter of constitution}).
    \item Pengawal demokrasi (\textit{the guardian of democracy}).
    \item Pelindung hak-hak konstitusional warga negara (\textit{the protector of citizen's constitutional rights}).
    \item Pelindung hak asasi manusia (\textit{the protector of human rights}).
\end{enumerate}


\subsection{7. Hak dan Kewajiban Warga Negara}
Hak berarti kepunyaan, kewenangan, kekuasaan untuk berbuat sesuatu (karena telah ditentukan oleh undang-undang, aturan, dan sebagainya). Adapun kewajiban artinya sesuatu yang harus dilaksanakan. Hak dan kewajiban merupakan sesuatu yang tidak dapat dipisahkan, tetapi sering terjadi pertentangan apabila pelaksanaannya tidak selaras, seimbang, dan proporsional.

Warga Negara Indonesia adalah orang-orang bangsa Indonesia asli dan orang-orang bangsa lain yang disahkan dengan undang-undang sebagai warga negara. Sejak menjadi warga negara maka melekatlah hak dan kewajiban sebagai Warga Negara Indonesia. Hak dan kewajiban yang dijamin oleh UUD NRI Tahun 1945 disebut sebagai hak dan kewajiban konstitusional. Adapun hak dan kewajiban yang diatur dalam peraturan perundang-undangan disebut sebagai hak dan kewajiban hukum.

\subsection{Sifat Pelaksanaan Hak dan Kewajiban Warga Negara}
\begin{enumerate}
    \item \textbf{Timbal Balik dengan Negara}: Hak warga negara merupakan kewajiban negara untuk memenuhinya, sedangkan kewajiban warga negara merupakan hak negara untuk memintanya.
    \item \textbf{Kausalitas atau Hubungan Sebab Akibat}: Seseorang memperoleh haknya karena dipenuhinya kewajiban yang dimilikinya.
    \item \textbf{Tidak Dapat Dipisahkan}: Dari kewajiban muncul hak dan dari hak muncul kewajiban.
\end{enumerate}

\subsection{Hak dan Kewajiban sebagaimana Dimuat dalam UUD NRI Tahun 1945}

\subsubsection{a. Terkait dengan Hukum dan Pemerintahan}
\begin{itemize}
    \item \textbf{Pasal 27 ayat (1)}: Segala warga negara bersamaan kedudukannya di dalam hukum dan pemerintahan dan wajib menjunjung hukum dan pemerintahan itu dengan tidak ada kecualinya.
    \item \textbf{Pasal 27 ayat (2)}: Tiap-tiap warga negara berhak atas pekerjaan dan penghidupan yang layak bagi kemanusiaan.
\end{itemize}

\subsubsection{b. Hak Asasi Manusia (HAM)}
\begin{itemize}
    \item \textbf{Pasal 28}: Kemerdekaan berserikat dan berkumpul, mengeluarkan pikiran dengan lisan dan tulisan dan sebagainya ditetapkan dengan undang-undang.
    \item \textbf{Pasal 28A}: Setiap orang berhak untuk hidup serta berhak mempertahankan hidup dan kehidupannya.
    \item \textbf{Pasal 28B}:
    \begin{enumerate}
        \item Setiap orang berhak membentuk keluarga dan melanjutkan keturunan melalui perkawinan yang sah.
        \item Setiap anak berhak atas kelangsungan hidup, tumbuh, dan berkembang serta berhak atas perlindungan dari kekerasan dan diskriminasi.
    \end{enumerate}
    \item \textbf{Pasal 28C}:
    \begin{enumerate}
        \item Setiap orang berhak mengembangkan diri melalui pemenuhan kebutuhan dasarnya, berhak mendapat pendidikan dan memperoleh manfaat dari ilmu pengetahuan dan teknologi, seni dan budaya, demi meningkatkan kualitas hidupnya dan demi kesejahteraan umat manusia.
        \item Setiap orang berhak untuk memajukan dirinya dalam memperjuangkan haknya secara kolektif untuk membangun masyarakat, bangsa, dan negaranya.
    \end{enumerate}
    \item \textbf{Pasal 28D}:
    \begin{enumerate}
        \item Setiap orang berhak atas pengakuan, jaminan, perlindungan, dan kepastian hukum yang adil serta perlakuan yang sama di hadapan hukum.
        \item Setiap orang berhak untuk bekerja serta mendapat imbalan dan perlakuan yang adil dan layak dalam hubungan kerja.
        \item Setiap warga negara berhak memperoleh kesempatan yang sama dalam pemerintahan.
        \item Setiap orang berhak atas status kewarganegaraan.
    \end{enumerate}
    \item \textbf{Pasal 28E}:
    \begin{enumerate}
        \item Setiap orang berhak memeluk agama dan beribadat menurut agamanya, memilih pendidikan dan pengajaran, memilih pekerjaan, memilih kewarganegaraan, memilih tempat tinggal di wilayah negara dan meninggalkannya, serta berhak kembali.
        \item Setiap orang berhak atas kebebasan meyakini kepercayaan, menyatakan pikiran dan sikap, sesuai dengan hati nuraninya.
        \item Setiap orang berhak atas kebebasan berserikat, berkumpul, dan mengeluarkan pendapat.
    \end{enumerate}
    \item \textbf{Pasal 28F}: Setiap orang berhak untuk berkomunikasi dan memperoleh informasi untuk mengembangkan pribadi dan lingkungan sosialnya, serta berhak untuk mencari, memperoleh, memiliki, menyimpan, mengolah, dan menyampaikan informasi dengan menggunakan segala jenis saluran yang tersedia.
    \item \textbf{Pasal 28G}:
    \begin{enumerate}
        \item Setiap orang berhak atas perlindungan diri pribadi, keluarga, kehormatan, martabat, dan harta benda yang di bawah kekuasaannya, serta berhak atas rasa aman dan perlindungan dari ancaman ketakutan untuk berbuat atau tidak berbuat sesuatu yang merupakan hak asasi.
        \item Setiap orang berhak untuk bebas dari penyiksaan atau perlakuan yang merendahkan derajat martabat manusia dan berhak memperoleh suaka politik dari negara lain.
    \end{enumerate}
    \item \textbf{Pasal 28H}:
    \begin{enumerate}
        \item Setiap orang berhak hidup sejahtera lahir dan batin, bertempat tinggal, dan mendapatkan lingkungan hidup yang baik dan sehat serta berhak memperoleh pelayanan kesehatan.
        \item Setiap orang berhak mendapat kemudahan dan perlakuan khusus untuk memperoleh kesempatan dan manfaat yang sama guna mencapai persamaan dan keadilan.
        \item Setiap orang berhak atas jaminan sosial yang memungkinkan pengembangan dirinya secara utuh sebagai manusia yang bermartabat.
        \item Setiap orang berhak mempunyai hak milik pribadi dan hak milik tersebut tidak boleh diambil alih secara sewenang-wenang oleh siapa pun.
    \end{enumerate}
    \item \textbf{Pasal 28I}:
    \begin{enumerate}
        \item Hak untuk hidup, hak untuk tidak disiksa, hak untuk kemerdekaan pikiran dan hati nurani, hak beragama, hak untuk tidak diperbudak, hak untuk diakui sebagai pribadi di hadapan hukum, dan hak untuk tidak dituntut atas dasar hukum yang berlaku surut adalah hak asasi manusia yang tidak dapat dikurangi dalam keadaan apa pun.
        \item Setiap orang bebas dari perlakuan yang bersifat diskriminatif atas dasar apa pun dan berhak mendapatkan perlindungan terhadap perlakuan yang bersifat diskriminatif itu.
        \item Identitas budaya dan hak masyarakat tradisional dihormati selaras dengan perkembangan zaman dan peradaban.
        \item Perlindungan, pemajuan, penegakan, dan pemenuhan hak asasi manusia adalah tanggung jawab negara, terutama pemerintah.
        \item Untuk menegakkan dan melindungi hak asasi manusia sesuai dengan prinsip negara hukum yang demokratis maka pelaksanaan hak asasi manusia dijamin, diatur, dan dituangkan dalam peraturan perundang-undangan.
    \end{enumerate}
    \item \textbf{Pasal 28J}:
    \begin{enumerate}
        \item Setiap orang wajib menghormati hak asasi manusia orang lain dalam tertib kehidupan bermasyarakat, berbangsa, dan bernegara.
        \item Dalam menjalankan hak dan kebebasannya, setiap orang wajib tunduk kepada pembatasan yang ditetapkan dengan undang-undang dengan maksud semata-mata untuk menjamin pengakuan serta penghormatan atas hak dan kebebasan orang lain dan untuk memenuhi tuntutan yang adil sesuai dengan pertimbangan moral, nilai-nilai agama, keamanan, dan ketertiban umum dalam suatu masyarakat demokratis.
    \end{enumerate}
\end{itemize}

\subsubsection{c. Jaminan Kebebasan Beragama atau Berkepercayaan}
\begin{itemize}
    \item \textbf{Pasal 29}:
    \begin{enumerate}
        \item Negara berdasar atas Ketuhanan Yang Maha Esa.
        \item Negara menjamin kemerdekaan tiap-tiap penduduk untuk memeluk agamanya masing-masing dan untuk beribadat menurut agamanya dan kepercayaannya itu.
    \end{enumerate}
\end{itemize}

\subsubsection{d. Bela Negara}
\begin{itemize}
    \item \textbf{Pasal 27 ayat (3)}: Setiap warga negara berhak dan wajib ikut serta dalam upaya pembelaan negara.
    \item \textbf{Pasal 30 ayat (1)}: Tiap-tiap warga negara berhak dan wajib ikut serta dalam usaha pertahanan dan keamanan negara.
\end{itemize}

\subsubsection{e. Pendidikan dan Kebudayaan}
\begin{itemize}
    \item \textbf{Pasal 31}:
    \begin{enumerate}
        \item Setiap warga negara berhak mendapat pendidikan.
        \item Setiap warga negara wajib mengikuti pendidikan dasar dan pemerintah wajib membiayainya.
        \item Pemerintah mengusahakan dan menyelenggarakan satu sistem pendidikan nasional, yang meningkatkan keimanan dan ketakwaan serta akhlak mulia dalam rangka mencerdaskan kehidupan bangsa, yang diatur dengan undang-undang.
        \item Negara memprioritaskan anggaran pendidikan sekurang-kurangnya dua puluh persen dari anggaran pendapatan dan belanja negara serta dari anggaran pendapatan dan belanja daerah untuk memenuhi kebutuhan penyelenggaraan pendidikan nasional.
        \item Pemerintah memajukan ilmu pengetahuan dan teknologi dengan menjunjung tinggi nilai-nilai agama dan persatuan bangsa untuk kemajuan peradaban serta kesejahteraan umat manusia.
    \end{enumerate}
    \item \textbf{Pasal 32}:
    \begin{enumerate}
        \item Negara memajukan kebudayaan nasional Indonesia di tengah peradaban dunia dengan menjamin kebebasan masyarakat dalam memelihara dan mengembangkan nilai-nilai budayanya.
        \item Negara menghormati dan memelihara bahasa daerah sebagai kekayaan budaya nasional.
    \end{enumerate}
\end{itemize}

\subsubsection{f. Perekonomian Nasional dan Kesejahteraan Sosial}
\begin{itemize}
    \item \textbf{Pasal 33}:
    \begin{enumerate}
        \item Perekonomian disusun sebagai usaha bersama berdasar atas asas kekeluargaan.
        \item Cabang-cabang produksi yang penting bagi negara dan yang menguasai hajat hidup orang banyak dikuasai oleh negara.
        \item Bumi dan air dan kekayaan alam yang terkandung di dalamnya dikuasai oleh negara dan dipergunakan untuk sebesar-besar kemakmuran rakyat.
        \item Perekonomian nasional diselenggarakan berdasar atas demokrasi ekonomi dengan prinsip kebersamaan, efisiensi berkeadilan, berkelanjutan, berwawasan lingkungan, kemandirian, serta dengan menjaga keseimbangan kemajuan dan kesatuan ekonomi nasional.
        \item Ketentuan lebih lanjut mengenai pelaksanaan pasal ini diatur dalam undang-undang.
    \end{enumerate}
    \item \textbf{Pasal 34}:
    \begin{enumerate}
        \item Fakir miskin dan anak-anak terlantar dipelihara oleh negara.
        \item Negara mengembangkan sistem jaminan sosial bagi seluruh rakyat dan memberdayakan masyarakat yang lemah dan tidak mampu sesuai dengan martabat kemanusiaan.
        \item Negara bertanggung jawab atas penyediaan fasilitas pelayanan kesehatan dan fasilitas pelayanan umum yang layak.
        \item Ketentuan lebih lanjut mengenai pelaksanaan pasal ini diatur dalam undang-undang.
    \end{enumerate}
\end{itemize}

\subsection{Contoh Hak dan Kewajiban Warga Negara Indonesia dalam Kehidupan Sehari-hari}

\subsubsection{Contoh Hak Warga Negara Indonesia}
\begin{enumerate}
    \item Memeluk dan menjalankan ibadah sesuai dengan agama dan kepercayaannya.
    \item Menyuarakan pendapat secara lisan maupun tulisan.
    \item Mendapatkan pendidikan dan pengajaran.
    \item Menikah dan berkeluarga.
    \item Mendapatkan penghidupan yang layak.
    \item Berorganisasi dan berkumpul.
    \item Memiliki suatu barang atau properti yang legal.
    \item Memiliki kedudukan yang sama di dalam hukum dan mendapat perlindungan hukum.
\end{enumerate}

\subsubsection{Contoh Kewajiban Warga Negara Indonesia}
\begin{enumerate}
    \item Membayar pajak dan retribusi yang ditetapkan pemerintah.
    \item Taat dan patuh pada peraturan hukum yang berlaku.
    \item Menghormati orang lain.
    \item Mengikuti pendidikan dasar.
    \item Ikut menjaga sarana dan fasilitas umum.
    \item Membuang sampah pada tempatnya.
    \item Patuh kepada pembatasan atas hak kebebasan.
\end{enumerate}